% ==============================================================
\section{Problem Definition}\label{sec:problem}
% ==============================================================

\subsection{Notation}

Let $\mathcal{I} = \{i_1, i_2, \ldots, i_m\}$ be a finite set of items,
and $D = [d_1, d_2, \ldots, d_N]$ be an ordered sequence of transactions
indexed by time.
Each transaction $d_t$ ($t = 1, \ldots, N$) is a subset of $\mathcal{I}$.
An itemset (pattern) $P \subseteq \mathcal{I}$ \emph{occurs} in transaction $d_t$ if $P \subseteq d_t$.

\begin{definition}[Global Support]\label{def:global_support}
The \emph{global support} of an itemset $P$ in $D$ is:
\begin{equation}
\text{gsupp}(P) = \frac{|\{t \mid 1 \leq t \leq N,\; P \subseteq d_t\}|}{N}
\label{eq:global_support}
\end{equation}
This measures the fraction of all transactions that contain $P$.
\end{definition}

\begin{definition}[Local Support (Windowed)]\label{def:local_support}
Given a window size $W \in \mathbb{Z}_{>0}$ and a window position $l \in \{1, \ldots, N - W + 1\}$,
the \emph{local support count} of $P$ at position $l$ is:
\begin{equation}
\text{lsup}(P, l, W) = |\{d_j \in D \mid l \leq j < l + W,\; P \subseteq d_j\}|
\label{eq:local_support}
\end{equation}
This counts the occurrences of $P$ within the sliding window $[l, l+W)$.
\end{definition}

\begin{definition}[Dense Interval]\label{def:dense_interval}
Given a density threshold $\theta \in \mathbb{Z}_{>0}$,
a \emph{dense interval} of $P$ is a maximal contiguous range $[s, e]$ of window positions such that:
\begin{equation}
\forall l \in [s, e]: \quad \text{lsup}(P, l, W) \geq \theta
\end{equation}
Let $\mathcal{D}(P, W, \theta)$ denote the set of all dense intervals of $P$.
\end{definition}

\subsection{Rare Dense Pattern}

We now introduce the central concept of this paper.

\begin{definition}[Global Rarity Condition]\label{def:global_rarity}
An itemset $P$ satisfies the \emph{global rarity condition} with respect to maximum support threshold $\sigma_{\max} \in (0, 1)$ if:
\begin{equation}
\text{gsupp}(P) < \sigma_{\max}
\label{eq:rarity}
\end{equation}
\end{definition}

\begin{definition}[Local Density Condition]\label{def:local_density}
An itemset $P$ satisfies the \emph{local density condition} with respect to window size $W$ and density threshold $\theta$ if:
\begin{equation}
\mathcal{D}(P, W, \theta) \neq \emptyset
\label{eq:density}
\end{equation}
That is, there exists at least one dense interval.
\end{definition}

\begin{definition}[Rare Dense Pattern (RDP)]\label{def:rdp}
An itemset $P$ is a \emph{Rare Dense Pattern} if it satisfies both:
\begin{enumerate}
\item The global rarity condition (Def.~\ref{def:global_rarity}): $\text{gsupp}(P) < \sigma_{\max}$
\item The local density condition (Def.~\ref{def:local_density}): $\mathcal{D}(P, W, \theta) \neq \emptyset$
\end{enumerate}
The set of all RDPs is denoted $\mathcal{R}(D, W, \theta, \sigma_{\max})$.
\end{definition}

\begin{remark}
The key tension is between the two conditions: $\sigma_{\max}$ enforces global rarity (the pattern must be \emph{infrequent} overall), while $\theta$ enforces local density (the pattern must be \emph{frequent} within some temporal window). The interplay $\theta / W > \sigma_{\max}$ ensures that a non-trivial class of patterns exists---those whose occurrences are temporally concentrated rather than uniformly distributed.
\end{remark}

\subsection{Two-Phase Mining}

\begin{definition}[Two-Phase Mining]\label{def:two_phase}
The \emph{Two-Phase Rare Dense Pattern Mining} algorithm operates as follows:
\begin{enumerate}
\item \textbf{Phase~1 (Dense Interval Discovery)}: Using a local density threshold $\theta$ and window size $W$, apply sliding-window Apriori to discover all itemsets $P$ with $\mathcal{D}(P, W, \theta) \neq \emptyset$. Denote this set $\mathcal{F}_{\text{local}}$.
\item \textbf{Phase~2 (Global Rarity Filtering)}: For each $P \in \mathcal{F}_{\text{local}}$, compute $\text{gsupp}(P)$ and retain only those with $\text{gsupp}(P) < \sigma_{\max}$:
\begin{equation}
\mathcal{R} = \{P \in \mathcal{F}_{\text{local}} \mid \text{gsupp}(P) < \sigma_{\max}\}
\end{equation}
\end{enumerate}
\end{definition}

\subsection{Properties and Theorems}

\begin{theorem}[Weak Anti-Monotonicity of Local Density]\label{thm:weak_antimon}
Let $P$ be an itemset satisfying the local density condition with parameters $(W, \theta)$.
Then for every subset $P' \subset P$, $P'$ also satisfies the local density condition
with the same parameters $(W, \theta)$.

Formally:
\begin{equation}
\mathcal{D}(P, W, \theta) \neq \emptyset \implies \forall P' \subset P: \mathcal{D}(P', W, \theta) \neq \emptyset
\end{equation}
\end{theorem}

\begin{proof}
Let $[s, e] \in \mathcal{D}(P, W, \theta)$ be a dense interval of $P$.
For any window position $l \in [s, e]$:
\[
\text{lsup}(P, l, W) = |\{j \mid l \leq j < l+W,\; P \subseteq d_j\}|
\]
Since $P' \subset P$, we have $P \subseteq d_j \implies P' \subseteq d_j$, and therefore:
\[
\text{lsup}(P', l, W) \geq \text{lsup}(P, l, W) \geq \theta
\]
Thus $[s, e]$ is also a dense interval for $P'$ (though $P'$ may have larger dense intervals that contain $[s,e]$).
Therefore $\mathcal{D}(P', W, \theta) \neq \emptyset$.
\end{proof}

\begin{remark}
We call this ``weak'' anti-monotonicity because the \emph{rarity} condition is \emph{monotone} (not anti-monotone): if $P$ is globally rare, supersets of $P$ are also rare.
Thus, for the \emph{combined} RDP condition, neither classical anti-monotonicity nor monotonicity holds in general.
However, the local density condition alone is anti-monotone, which enables Apriori-style pruning in Phase~1.
\end{remark}

\begin{theorem}[Monotonicity of Global Rarity]\label{thm:rarity_monotone}
If $\text{gsupp}(P) < \sigma_{\max}$, then for every superset $Q \supset P$:
\begin{equation}
\text{gsupp}(Q) \leq \text{gsupp}(P) < \sigma_{\max}
\end{equation}
\end{theorem}

\begin{proof}
Immediate from $Q \supset P$: $P \subseteq d_t$ does not imply $Q \subseteq d_t$,
but $Q \subseteq d_t \implies P \subseteq d_t$.
Thus $\{t \mid Q \subseteq d_t\} \subseteq \{t \mid P \subseteq d_t\}$,
giving $\text{gsupp}(Q) \leq \text{gsupp}(P)$.
\end{proof}

\begin{theorem}[Completeness of Two-Phase Mining]\label{thm:completeness}
The Two-Phase Mining algorithm (Def.~\ref{def:two_phase}) discovers exactly the set $\mathcal{R}(D, W, \theta, \sigma_{\max})$ of all Rare Dense Patterns.
\end{theorem}

\begin{proof}
We prove both directions:

\textbf{Soundness ($\mathcal{R}_{\text{output}} \subseteq \mathcal{R}$):}
Every pattern $P$ output by the algorithm satisfies both conditions by construction:
Phase~1 ensures $\mathcal{D}(P, W, \theta) \neq \emptyset$, and
Phase~2 ensures $\text{gsupp}(P) < \sigma_{\max}$.

\textbf{Completeness ($\mathcal{R} \subseteq \mathcal{R}_{\text{output}}$):}
Let $P \in \mathcal{R}$ be any RDP. We must show that Phase~1 discovers $P$.
Phase~1 uses Apriori with pruning based on local density.
By Theorem~\ref{thm:weak_antimon}, if $P$ has a dense interval, then every subset $P' \subset P$ also has a dense interval.
Therefore no subset of $P$ is pruned by Apriori's anti-monotone pruning, and $P$ will be generated as a candidate.
Since $P$ itself has a dense interval, it passes the Phase~1 filter.
Phase~2 then checks global rarity; since $P \in \mathcal{R}$, it satisfies $\text{gsupp}(P) < \sigma_{\max}$ and is retained.
\end{proof}

\begin{theorem}[Pruning Correctness]\label{thm:pruning}
In Phase~1, if an itemset $P'$ has $\mathcal{D}(P', W, \theta) = \emptyset$
(no dense interval), then every superset $Q \supset P'$ also has
$\mathcal{D}(Q, W, \theta) = \emptyset$, and can be safely pruned.
\end{theorem}

\begin{proof}
This is the contrapositive of Theorem~\ref{thm:weak_antimon}:
if $Q$ had a dense interval, then $P' \subset Q$ would also have a dense interval.
Since $P'$ does not, neither does $Q$.
\end{proof}

\subsection{Complexity Analysis}

\begin{theorem}[Time Complexity]\label{thm:complexity}
Let $m = |\mathcal{I}|$ be the number of distinct items,
$N$ the number of transactions, $k_{\max}$ the maximum itemset size,
and $C_k$ the number of candidate $k$-itemsets at level $k$.
The time complexity of Two-Phase Mining is:
\begin{equation}
O\Bigl(\sum_{k=1}^{k_{\max}} C_k \cdot N \log N\Bigr)
\end{equation}
where the $N \log N$ factor per candidate arises from binary search in the dense interval computation.
\end{theorem}

\begin{proof}
For each candidate itemset at level $k$:
\begin{enumerate}
\item Computing co-occurrence timestamps via sorted list intersection: $O(N)$
\item Computing dense intervals via sliding window with binary search: $O(N \log N)$
\end{enumerate}
Phase~2 computes global support in $O(N)$ per pattern and is dominated by Phase~1.

The key efficiency gain over na\"ive enumeration comes from Apriori pruning:
by Theorem~\ref{thm:pruning}, itemsets without dense intervals and their supersets are pruned,
reducing $C_k$ significantly in practice.
\end{proof}

\begin{remark}[Comparison with Standard Apriori]
Standard Apriori with a global minimum support threshold $\sigma_{\min}$ prunes \emph{infrequent} itemsets.
Our approach prunes itemsets without dense intervals instead.
Since we target globally rare patterns, standard Apriori would prune exactly the patterns we seek.
The local density condition provides an alternative anti-monotone criterion that enables efficient search while preserving the rare patterns of interest.
\end{remark}
